% frontmatter/Abstract.tex — dissertation abstract.
% Plain Ch.4 register per outlines/style_guide.md. Regenerated 2026-06-09 via the
% regenerate-not-patch workflow (style_guide.md §6): re-authored from a bare-claims
% skeleton conditioned on Ch.4 rhythm spans, gated with voice_meter.py, and varied
% asymmetrically against the Ch.1 thesis statement so the shared n-grams clear.
% Must still fit ONE PAGE — keep the \linespread line; do not lengthen. Framing locks:
% representation-structure thesis, BOTH lenses; arc failure -> recovery -> detection;
% "guaranteed failure" Ch.4-only; Ch.5 diagnostic+recovery; Ch.6 = parity / "matches
% or, in the mean, outperforms", never "beats"/"superior". Bookend sentence "When the
% representation keeps it, a detector can read it out." stays verbatim (sanctioned).
\begin{abstract}
% Tighten body leading slightly so the abstract fits on one page (the global
% \linespread{1.25} pushed the final sentence onto a second page). Scoped to this
% environment's group, so it does not affect any other page.
\linespread{1.12}\selectfont
Models deployed outside controlled settings encounter inputs unlike their training
data and may generate outputs unsupported by what they learned. A reliable system
should recognize these cases. Out-of-distribution detection and hallucination
detection are usually studied separately, but both depend on the same condition:
the model's representation must retain the information that distinguishes an
anomalous case from an ordinary one. If training never captures that information,
or later suppresses it, no downstream detector can recover it.

We study this condition from two complementary perspectives. Mutual information
measures how much information about the relevant distinction remains in a
representation, while feature-space geometry shows whether the directions needed to
detect a shift retain sufficient variation. The first study proves that unlabeled
detection must fail when the self-supervised or unsupervised objective is independent
of the label-relevant features. It also introduces the Adjacent OOD benchmark to
evaluate this high-overlap regime. The second study identifies domain-sensitivity
collapse, in which single-domain training suppresses shift-sensitive directions, and
introduces Teacher-Guided Training to restore them without adding inference cost. The
third study develops a hallucination detector from cross-layer activations. Its
contrastive probe reaches parity with the strongest engineered probe and outperforms
it in the mean, consistent with the study's feasibility argument.

The unifying representation-structure framework is the central contribution. The
dissertation establishes a conditional limit on unlabeled OOD detection and
demonstrates interventions when relevant structure can be restored or read out; it
does not claim these methods are optimal.
\end{abstract}
